\section{Introduction}\label{sec:intro}

Large language model (LLM) based agents are increasingly deployed as
autonomous decision-makers in financial markets, where they execute trading
strategies with minimal human oversight~\cite{xiao2024tradingagents,
liu2025finrs}. In such settings, the agent optimizes a proxy reward that
approximates, but does not perfectly capture, the principal's true objective.
When the proxy reward is susceptible to perturbation---whether through
adversarial manipulation, distributional shift, or specification error---the
agent's behavior may deviate from the intended objective, a phenomenon broadly
termed \emph{misalignment}. Given the economic stakes of automated trading, a
quantitative understanding of when misalignment becomes behaviorally
significant is essential for responsible deployment.

The theoretical foundations of misalignment have advanced considerably.
Skalse et al.~\cite{skalse2022defining} proved that reward hacking is
generically unavoidable: over any open set of stochastic policies, every
non-trivial proxy reward is hackable. Wang and Huang~\cite{wang2026reward}
showed that distortion is inevitable under finite evaluation, and conjectured
the existence of a phase transition between gaming and degrading regimes.
On the empirical side, Debenedetti et al.~\cite{debenedetti2023scaling}
established a power-law relationship between compute and adversarial
robustness with a remarkably small exponent ($\alpha \approx 0.01$),
while Nathanson et al.~\cite{nathanson2026scaling} demonstrated logarithmic
scaling of adversarial harm with model size ratios.

Despite these advances, no existing framework provides a quantitative answer
to the following question: \emph{given a trading agent with specific policy
parameters operating in a market with given volatility and decision horizon,
what is the minimum perturbation intensity that induces statistically
significant misalignment?} The impossibility results of
\cite{skalse2022defining} guarantee that misalignment can occur but do not
predict the perturbation magnitude required. The distortion analysis of
\cite{wang2026reward} characterizes the direction of misalignment but not
the threshold intensity. The empirical scaling laws of
\cite{debenedetti2023scaling, nathanson2026scaling} lack theoretical
derivation and do not connect to market-specific parameters.

In this paper, we develop a rigorous mathematical framework for
\emph{misalignment sensitivity thresholds} in the context of
softmax-parameterized trading agents operating in finite Markov decision
processes (MDPs). Our contributions are as follows.

\begin{itemize}[leftmargin=*,itemsep=2pt,topsep=4pt]
\item We introduce the misalignment sensitivity threshold $\tau(\delta)$,
  defined as the infimum perturbation intensity $\varepsilon$ for which the
  KL divergence between perturbed and baseline policies exceeds a
  significance level $\delta$ (Definition~\ref{def:threshold}).
\item We prove that $\tau(\delta)$ exists and is finite for any finite MDP
  with softmax policies (Theorem~\ref{thm:existence}), and characterize it
  via the Fisher information of the policy parameterization:
  $\tau(\delta) = \sqrt{2\delta/\lVert G \rVert_F^2} + O(\delta)$
  (Theorem~\ref{thm:characterization}).
\item We establish the temperature scaling law $\tau \propto T$
  (Corollary~\ref{cor:temperature}), identifying the softmax temperature as
  the dominant determinant of misalignment robustness.
\item We prove that market volatility $\sigma$ and decision horizon $h$ are
  asymptotically negligible, with $\tau(\sigma) = \tau_0 + O(\sigma^{-2})$
  and $\tau(h) = \tau_0 + O(h^{-1})$ (Propositions~\ref{prop:volatility}
  and~\ref{prop:horizon}), refuting the initially hypothesized power-law
  dependence $\tau \propto \sigma^{-\alpha} h^{-\beta}$.
\item We verify all theoretical predictions computationally across more than
  one hundred MDP configurations, confirming the linear temperature scaling
  with $R^2 = 0.999$ and the negligible market-parameter exponents
  $|\alpha|, |\beta| < 0.02$.
\end{itemize}

The proofs combine perturbation analysis of softmax policies, second-order
Taylor expansion of KL divergence, and the standard softmax gradient identity.
The key technical insight is that the Fisher information of a softmax
distribution concentrates the sensitivity of the policy in a single
parameter---the temperature $T$---that governs the ``peakedness'' of the
decision distribution. Market parameters affect the Q-values but not the
mechanism by which perturbations propagate through the softmax, leading to
their asymptotic irrelevance.

The paper is organized as follows. Section~\ref{sec:prelim} establishes
notation and defines the trading MDP, intervention operators, and the
misalignment threshold. Section~\ref{sec:main} states and proves the main
results. Section~\ref{sec:discussion} discusses implications, connections to
existing work, and open questions. Section~\ref{sec:conclusion} summarizes
our contributions.
